\begin{figure}[!ht]
  \centering
  \includegraphics[width=\linewidth]{assets/wiggle_framework.pdf}
  \caption{\textbf{The Wiggle Framework.} Every trajectory is anchored to L0, the first valid verdict produced by the temperature-zero, no-pressure applied. 
  \emph{Mechanical Consistency} re-attempts the query in a semantically invariant manner (basic repetition, trivial prompt perturbation, positional consistency). 
  \emph{Single-turn Conviction} measures the response to one scripted challenge (L1--L4).
  \emph{Multi-turn Persistence} sustains L1--L4 for 10 turns and adds two expressly multi-turn protocols: cycling through L1--L4 (L5) 
  and a separate LLM generating each challenge adaptively from the conversation so far (L6). 
  \texttt{[opposite]} is the verdict opposite to L0 and \texttt{[argument]} is a model-generated argument (Appendix~\ref{sec:appendix-pressure-gen}). 
  Full L6 protocol in Appendix~\ref{sec:appendix-l6-protocol}.}
  \label{fig:framework}
\end{figure}

\section{Introduction}
\label{sec:introduction}

LLM judges sit at increasingly consequential decision points across the model development stack. 
They score outputs in benchmarks, classify content in production, and stand in for human judgment 
in the loops that train, grade, and refine frontier models. The standard validation workflow is to 
curate a golden set of expert-vetted examples, verify that the judge's verdicts align reasonably with those labels, 
and deploy if accuracy is sufficient~\citep{collot2025balanced}. 
This establishes whether a judge is correct \emph{on average}, but it says much less about whether it is 
\emph{stable}. Would the verdict survive if the judge were asked again, challenged, or pressed repeatedly? 
Behind every verdict sits a hidden distribution over how much conviction the judge actually holds.

To pinpoint how confident a judge is, several strategies have been proposed:

\begin{itemize}
  \item \textbf{Just ask.} Ask the LLM to produce a self-reported confidence score (e.g., ``how confident are you, 0–100\%?'').
  Models have been shown to be badly miscalibrated in the overconfident direction~\citep{wei2024simpleqa}.
  \item \textbf{Observe consistency over many repetitions.} Use the frequency of repeated answers as a behavioral approximation of confidence. 
  This adds substantial inference cost while inheriting the same overconfidence problem~\citep{wei2024simpleqa}.
  \item \textbf{Inspect verdict-token log probabilities.} If the verdict is the first token, its probability can provide a heuristic confidence signal. 
  Reasoning models deliberate before answering, and many frontier APIs no longer expose raw token probabilities, in part because such outputs can
  enable model-extraction attacks and leak proprietary model information~\citep{carlini2024stealing,finlayson2024logits}.
\end{itemize}

A largely separate literature studies \emph{sycophancy} and \emph{persuadability}: the tendency of LLMs to fold under conversational pressure,
plausibly a side-effect of preference optimization ~\citep{perez2022discovering,sharma2024sycophancy,laban2023flipflop}. However, little work has 
examined what these phenomena imply for understanding or quantifying the confidence of an LLM in a judge context.

We introduce the \textbf{Wiggle Framework}, a unified stress test for epistemic stability in LLM judges. 
It decomposes judge confidence into three behaviorally grounded dimensions: \emph{Mechanical Consistency} (stability under re-prompting and semantically invariant 
prompt variation), \emph{Single-turn Conviction} (stability under a single substantive challenge), and \emph{Multi-turn Persistence} (stability under sustained or adaptive
pressure). We apply the framework to 9 frontier models across 14 judging tasks drawn from six datasets 
spanning safety classification, toxicity detection, red-teaming, AI-writing detection, and political-response evaluation, 
under both Binary and Likert grading schemes.

All models tested as judges wiggle at substantial rates across all datasets and grading schemes (Figure~\ref{fig:hero}). The
fact that LLMs change their minds under pressure is not new, but what is more surprising is the
\emph{structure} of the flips: Binary and Likert scales produce opposite directional tendencies on the same items, 
and when a judge does flip, the flip is far more often corruptive than corrective with respect to ground truth.
Section \S\ref{sec:discussion} discusses several findings, including correlations in susceptibility across pressure 
levels, the limited transferability of wiggle-rate profiles across datasets and models, baseline agreement within a model 
jury as the strongest inexpensive predictor of item-level epistemic instability, and the implications of 
predominantly corruptive pressure for deploying judges in self-governing agentic systems.
